\chapter{Huisken monotonicity and entropy}\label{sec:mcf-monotonicity}
The Gaussian integral measures area near a spacetime point at the parabolic scale. Its dissipation identifies self-similar models, while taking the supremum over centers and scales gives entropy. We first prove the identities for smooth closed Euclidean flows, so every integration by parts has no boundary term. Noncompact and weak formulations need additional hypotheses and are stated separately. The main source is \cite{Hui90}, Section 3, with the normalization converted to that of Definition \ref{def:self-shrinker}.

\begin{proposition}[Transport of a test function along a closed flow]\label{prop:mcf-test-transport}
Let $\Sigma_t$ be a smooth closed mean curvature flow in $\mathbb R^{n+1}$, and let $\phi(x,t)$ be smooth in an ambient neighborhood of its spacetime track. Then
\[
 \frac{d}{dt}\int_{\Sigma_t}\phi\,d\mu_t
 =\int_{\Sigma_t}\left(\partial_t\phi+\nabla\phi\cdot\vec H-\phi|\vec H|^2\right)d\mu_t.
\]
Here $\partial_t\phi$ is the ambient partial derivative, with the spatial point held fixed.
\end{proposition}
\begin{proof}
Choose a normal parametrization. The chain rule gives $\frac{d}{dt}\phi(F(p,t),t)=\partial_t\phi+\nabla\phi\cdot\vec H$, and Proposition \ref{prop:normal-variation-formulas} gives $\partial_t d\mu=-|\vec H|^2d\mu$. Adding the two contributions proves the formula. A tangential velocity adds a tangential divergence, whose integral vanishes on a closed hypersurface. Compare \cite{WhiMCF15}, Sections 6 and 9.
\end{proof}

\begin{definition}[Backward Gaussian kernel]\label{def:mcf-backward-gaussian}
For a spacetime center $(x_0,t_0)$ and $t<t_0$, define
\[
 \rho_{x_0,t_0}(x,t)=(4\pi(t_0-t))^{-n/2}
 \exp\left(-\frac{|x-x_0|^2}{4(t_0-t)}\right).
\]
The normalization uses the dimension $n$ of the evolving hypersurface, not the ambient dimension $n+1$. In particular its integral over an affine $n$-plane through $x_0$ is one. See \cite{Hui90}, p.~288.
\end{definition}

\begin{theorem}[Huisken's monotonicity identity for closed smooth flows]\label{thm:huisken-monotonicity}
For a smooth closed Euclidean mean curvature flow and a fixed $(x_0,t_0)$, at every flow time $t<t_0$,
\[
 \frac{d}{dt}\int_{\Sigma_t}\rho_{x_0,t_0}\,d\mu_t
 =-\int_{\Sigma_t}\left|\vec H+
 \frac{(x-x_0)^\perp}{2(t_0-t)}\right|^2\rho_{x_0,t_0}\,d\mu_t.
\]
This is \cite{Hui90}, Theorem 3.1, pp.~288--289. The square uses the vector mean curvature, so the sign agrees with $\vec H=-H\nu$.
\end{theorem}
\begin{proof}
Write $y=x-x_0$, $\tau=t_0-t$, and $\rho=\rho_{x_0,t_0}$. Direct differentiation gives
\[
 \partial_t\rho=\rho\left(\frac n{2\tau}-\frac{|y|^2}{4\tau^2}\right),
 \qquad \nabla\rho=-\frac{y}{2\tau}\rho.
\]
The integrand in Proposition \ref{prop:mcf-test-transport} is consequently
\[
 \rho\left(\frac n{2\tau}-\frac{|y|^2}{4\tau^2}
 -\frac{\langle y,\vec H\rangle}{2\tau}-|\vec H|^2\right).
\]
Since $\operatorname{div}_{\Sigma_t}y^T=n+\langle y,\vec H\rangle$, the product rule gives
\[
 \operatorname{div}_{\Sigma_t}\left(\rho\frac{y^T}{2\tau}\right)
 =\rho\left(\frac n{2\tau}+\frac{\langle y,\vec H\rangle}{2\tau}
 -\frac{|y^T|^2}{4\tau^2}\right).
\]
Subtract this divergence from the previous integrand. The remaining expression is $-\rho|\vec H+y^\perp/(2\tau)|^2$. The divergence integrates to zero because the time slice is compact and without boundary.
\end{proof}

\begin{corollary}[Gaussian density at an endpoint]\label{cor:mcf-gaussian-density-limit}
Suppose a smooth closed flow is defined on $[a,t_0)$. For every $x_0\in\mathbb R^{n+1}$ the finite limit
\[
 \Theta(\mathcal M,x_0,t_0)=
 \lim_{t\uparrow t_0}\int_{\Sigma_t}\rho_{x_0,t_0}\,d\mu_t
\]
exists and is nonnegative. Equivalently, the Gaussian ratio at scale $r$ is the integral on $\Sigma_{t_0-r^2}$ with variance $r^2$, and these ratios are nondecreasing as $r$ increases within the available time interval.
\end{corollary}
\begin{proof}
Theorem \ref{thm:huisken-monotonicity} makes the integral nonincreasing in time and bounded below by zero. At any fixed earlier time it is finite because the time slice is compact. This yields a finite one-sided limit even when the flow is singular at $t_0$. Replacing $t$ by $t_0-r^2$ reverses the order of the parameter. This is the smooth-flow density construction used in \cite{WhiMCF15}, Section 10.2.
\end{proof}

\begin{proposition}[Gaussian density at a smooth embedded point]\label{prop:mcf-regular-density}
If a closed embedded flow is smooth through time $t_0$ and $x_0\in\Sigma_{t_0}$, then $\Theta(\mathcal M,x_0,t_0)=1$. If $x_0\notin\Sigma_{t_0}$, the density is zero. These assertions concern a single embedded sheet with unit multiplicity.
\end{proposition}
\begin{proof}
Near a point on $\Sigma_{t_0}$, use a smooth graphical representation over its tangent plane. At time $t_0-r^2$, dilating by $r^{-1}$ about $x_0$ makes the graph converge smoothly on each bounded set to that plane: the spatial graph error is quadratic and the time displacement is $O(r^2)$. Its Gaussian integral converges to the plane integral, one. On the fixed graphical neighborhood the rescaled area elements are uniformly bounded and the Gaussian tails are uniformly integrable. Outside that neighborhood a fixed positive distance from $x_0$ gives an exponentially small term, since the total areas stay bounded near $t_0$. If $x_0$ is off the final slice, compactness and smooth dependence in time give a positive distance from all nearby slices, so the entire integral tends to zero. Compare \cite{WhiMCF15}, Section 10.2.
\end{proof}

\begin{definition}[Gaussian area at arbitrary center and scale, and entropy]\label{def:mcf-entropy}
For a Radon measure $\mu$ on $\mathbb R^{n+1}$, put
\[
 F_{x_0,s}(\mu)=(4\pi s)^{-n/2}\int e^{-|x-x_0|^2/(4s)}\,d\mu(x),
 \qquad s>0,
\]
\[
 \lambda(\mu)=\sup_{x_0\in\mathbb R^{n+1},\ s>0}F_{x_0,s}(\mu).
\]
The value may be infinite. For a hypersurface, use its area measure; for an integral varifold, use its weight measure, including multiplicity. The functional in Proposition \ref{prop:gaussian-first-variation} is $F_{0,1}$. Sources: \cite{WhiMCF15}, Section 10, and \cite{CCMS20}, Section 2.
\end{definition}

\begin{proposition}[Entropy is invariant under translations and dilations]\label{prop:mcf-entropy-invariance}
Let $D(x)=b+ax$, $a>0$, and define the geometric dilation of an $n$-dimensional measure by $\mu'=a^nD_\#\mu$. Then
\[
 F_{x_0,s}(\mu')=F_{(x_0-b)/a,s/a^2}(\mu),
 \qquad \lambda(\mu')=\lambda(\mu).
\]
Thus the measure scaling includes the tangential area factor, as in Definition \ref{def:varifold-dilation}.
\end{proposition}
\begin{proof}
Substitute $x=b+ay$ into the defining integral. The factor $a^n$ combines with $(4\pi s)^{-n/2}$ to give $(4\pi s/a^2)^{-n/2}$. The change $(x_0,s)\mapsto((x_0-b)/a,s/a^2)$ is a bijection on the set of centers and positive scales, so taking suprema proves the entropy identity.
\end{proof}

\begin{proposition}[Entropy is nonincreasing along a closed smooth flow]\label{prop:mcf-entropy-monotone}
If $t_1<t_2$ are smooth times of a closed Euclidean mean curvature flow, then
\[
 \lambda(\Sigma_{t_2})\leq\lambda(\Sigma_{t_1}).
\]
The supremum in the entropy definition need not be attained for this argument.
\end{proposition}
\begin{proof}
Fix $x_0$ and $s>0$ and apply Theorem \ref{thm:huisken-monotonicity} with terminal time $t_2+s$. It gives
\[
 F_{x_0,s}(\Sigma_{t_2})\leq
 F_{x_0,s+t_2-t_1}(\Sigma_{t_1})\leq\lambda(\Sigma_{t_1}).
\]
Take the supremum over $x_0,s$. This is the entropy consequence discussed in \cite{WhiMCF15}, Section 10.
\end{proof}

\begin{proposition}[Entropy and uniform Euclidean area growth]\label{prop:mcf-entropy-area-growth}
If $\lambda(\mu)\leq\Lambda$, then for every $x\in\mathbb R^{n+1}$ and $r>0$,
\[
 \mu(B_r(x))\leq e^{1/4}(4\pi)^{n/2}\Lambda r^n.
\]
Conversely, if $\mu(B_r(x))\leq Kr^n$ for every center and radius, then $\lambda(\mu)\leq C(n)K$. The uniformity in the center is part of the hypothesis.
\end{proposition}
\begin{proof}
In $B_r(x)$ the exponential in $F_{x,r^2}$ is at least $e^{-1/4}$, proving the first inequality. Conversely fix a scale $s=r^2$. Bound the contribution of $B_r(x)$ by $(4\pi)^{-n/2}K$. On the annulus $2^jr\leq|y-x|<2^{j+1}r$, the exponential is at most $e^{-4^j/4}$ and the mass is at most $K(2^{j+1}r)^n$. Thus the whole integral is bounded by
\[
 (4\pi)^{-n/2}K\left(1+\sum_{j=0}^\infty2^{(j+1)n}e^{-4^j/4}\right).
\]
The sum converges independently of $x,r$. Compare the bounded-area-ratio class used in \cite{WhiMCF15}, Section 10.1.
\end{proof}

\begin{example}[A plane records its multiplicity in entropy]\label{ex:mcf-plane-entropy}
For an affine $n$-plane $P$ with integer multiplicity $q\geq1$,
\[
 F_{x_0,s}(q\mathcal H^n\llcorner P)
 =q\exp\left(-\frac{\operatorname{dist}(x_0,P)^2}{4s}\right),
 \qquad \lambda(q|P|)=q.
\]
To check this, split $x-x_0$ into orthogonal tangential and normal components and integrate the tangential Gaussian. The supremum is attained by choosing $x_0\in P$. Small support excess therefore does not make the entropy close to one when multiplicity is large; compare Example \ref{ex:multiple-plane-excess}.
\end{example}

\chapter{Rescaled flow and self-similar models}\label{sec:mcf-rescaled-models}
Logarithmic time converts a finite-time singularity into a long-time problem. This is a change of variables before it is a compactness theorem. We retain the scale-one shrinker convention $\vec H+x^\perp/2=0$ throughout.

\begin{proposition}[Logarithmic rescaling at a finite time]\label{prop:mcf-log-rescaling}
Fix a point $x_0$ and a time $T$ later than the times under consideration. Set
\[
 \tau=-\log(T-t),\qquad
 \widetilde F(p,\tau)=(T-t)^{-1/2}(F(p,t)-x_0).
\]
Then the normal velocity of the rescaled hypersurfaces is
\[
 (\partial_\tau\widetilde F)^\perp=
 \vec H_{\widetilde\Sigma_\tau}+\tfrac12\widetilde F^\perp.
\]
If the original flow is Type I at $T$, its rescaled curvature is uniformly bounded. This last assertion supplies a curvature bound, not a claim of convergence or uniqueness of a limiting shrinker.
\end{proposition}
\begin{proof}
Since $dt/d\tau=T-t$ and $(T-t)^{-1/2}=e^{\tau/2}$, differentiation yields
$\partial_\tau\widetilde F=\widetilde F/2+(T-t)^{1/2}\vec H$.
The second term is the mean curvature vector of the dilated hypersurface. Project normally to get the equation. Proposition \ref{prop:mcf-parabolic-rescaling} gives $|\widetilde A|^2=(T-t)|A|^2$, proving the curvature assertion. Huisken \cite{Hui90}, Section 2, uses instead $(2(T-t))^{-1/2}$ and $-\frac12\log(T-t)$; his rescaled equation has coefficient one instead of $1/2$.
\end{proof}

\begin{proposition}[Gaussian area dissipates along rescaled flow]\label{prop:mcf-rescaled-dissipation}
For a smooth closed rescaled flow satisfying Proposition \ref{prop:mcf-log-rescaling},
\[
 \frac{d}{d\tau}F_{0,1}(\widetilde\Sigma_\tau)
 =-(4\pi)^{-n/2}\int_{\widetilde\Sigma_\tau}
 \left|\vec H+\tfrac12x^\perp\right|^2e^{-|x|^2/4}\,d\mu.
\]
Consequently, if the rescaled flow exists for all $\tau\geq\tau_0$, then the time integral of the nonnegative expression on the right, with its minus sign removed, is finite.
\end{proposition}
\begin{proof}
Insert the rescaled normal velocity in Proposition \ref{prop:gaussian-first-variation}. Tangential velocity contributes only a divergence. Integrate in $\tau$ and use $F_{0,1}\geq0$. This is \cite{Hui90}, Corollary 3.2, after converting the Gaussian normalization.
\end{proof}

\begin{corollary}[Times with small Gaussian dissipation]\label{cor:mcf-small-dissipation-times}
Under the long-time hypothesis in Proposition \ref{prop:mcf-rescaled-dissipation}, there are $\tau_j\to\infty$ such that
\[
 \int_{\widetilde\Sigma_{\tau_j}}
 \left|\vec H+\tfrac12x^\perp\right|^2e^{-|x|^2/4}\,d\mu\longrightarrow0.
\]
If, in addition, these slices converge smoothly with multiplicity one on compact subsets to a smooth hypersurface $\Sigma_\infty$, then $\Sigma_\infty$ satisfies the self-shrinker equation. The compactness assumption is separate from the dissipation conclusion.
\end{corollary}
\begin{proof}
A nonnegative integrable function on $[\tau_0,\infty)$ has a sequence of times tending to infinity at which it tends to zero. For the second assertion, multiply the integrand by a nonnegative compactly supported ambient cutoff. Smooth convergence permits passage to the limit in this localized integral, which is bounded above by the full dissipation. The limit integral is zero for every cutoff; continuity and the positive Gaussian weight force $\vec H+x^\perp/2=0$ on $\Sigma_\infty$. No assertion about mass escaping to infinity or uniqueness follows from this local argument. Compare \cite{WhiMCF15}, Section 6.
\end{proof}

\begin{proposition}[Equality in monotonicity gives self-similar motion]\label{prop:mcf-monotonicity-equality}
Suppose a smooth closed flow has constant Gaussian integral centered at $(x_0,T)$ on a time interval $(a,b)\subset(-\infty,T)$. Then throughout that interval
\[
 \vec H+\frac{(x-x_0)^\perp}{2(T-t)}=0.
\]
Equivalently, the images $(T-t)^{-1/2}(\Sigma_t-x_0)$ are independent of time and form a self-shrinker in the project's scale-one normalization.
\end{proposition}
\begin{proof}
Theorem \ref{thm:huisken-monotonicity} gives zero integral of a continuous nonnegative function at each time, and hence the displayed pointwise equation. By Proposition \ref{prop:mcf-log-rescaling}, the rescaled velocity is tangent everywhere. A time-dependent tangential vector field on the compact parameter manifold can be removed by reparametrization, so the rescaled images are constant. The equation then reads $\vec H+x^\perp/2=0$. Conversely that equation makes the monotonicity integrand zero.
\end{proof}

\begin{definition}[Translators and self-expanders]\label{def:mcf-other-solitons}
A smooth hypersurface is a translator with velocity $v\in\mathbb R^{n+1}$ if $\vec H=v^\perp$. Its translates $\Sigma_t=\Sigma+tv$ have normal velocity $\vec H$. A scale-one self-expander satisfies $\vec H-x^\perp/2=0$; its expanding family is $\Sigma_t=\sqrt t\,\Sigma$, $t>0$. These are exact solution classes, not automatic classifications of Type II singularities. Compare the translating models discussed in \cite{HP99}, Section 4, and \cite{WhiMCF15}, Section 9.3.
\end{definition}

\begin{example}[The grim reaper and its cylindrical products]\label{ex:mcf-grim-reaper}
The graphs
\[
 u(x,t)=-\log(\cos x)+t,\qquad -\pi/2<x<\pi/2,
\]
give a translating curve-shortening flow. Indeed $u_t=1$, $u_x=\tan x$, and $u_{xx}=\sec^2x$, so $u_t=u_{xx}/(1+u_x^2)$. Taking the product with $\mathbb R^{n-1}$ gives a translating hypersurface flow in $\mathbb R^{n+1}$ with the same scalar graphical equation. Its slope is unbounded near the boundary of the graph domain, even though the complete translating hypersurface has no boundary there. The intrinsic length to either end diverges because $\sqrt{1+u_x^2}=1/\cos x$. Compare the translator discussion in \cite{HP99}, Section 4.
\end{example}
